\documentclass[12pt]{scrartcl}
\usepackage[hagavi]{azzam}

% Fungsi Bulat - Soal Latihan
% 50 soal: 20 Banyak Faktor, 11 Jumlah Faktor, 4 Totient, 9 Campuran, 6 Tantangan.
% Sumber tiap soal dimunculkan di PDF sebagai \item (Sumber); komentar menyimpan detail lengkap.
% Soal pilihan ganda asli disajikan tanpa opsi jawaban sebagai soal isian.
% Versi gabungan Materi+Latihan tersedia sebagai Fungsi-Aritmetika_Modul.tex.

\title{Fungsi Aritmetika - Soal Latihan}
\date{15 Oktober 2026}

\begin{document}
\maketitle

\renewcommand*\contentsname{Daftar Isi}
\tableofcontents

\newpage

\section{Banyak Faktor}

\begin{enumerate}

	% Source: WMI 2019 Grade 5-6 (diterjemahkan)
	% Jawaban: 7
	\item (WMI 2019)
	Ada berapa bilangan asli antara $300$ sampai $600$ yang banyak faktor
	positifnya ganjil?

	% Source: AMO 2019 Grade 8 No. 18 (diterjemahkan)
	% Jawaban: 7
	\item (AMO 2019)
	Ada berapa bilangan $3$ digit yang tepat memiliki $3$ faktor positif?

	% Source: OSK SMA 2015 No. 1 (Isian Singkat)
	% Jawaban: 8
	\item (OSK 2015)
	Banyaknya faktor bulat positif dari 2015 adalah \dots

	% Source: OSK SMA 2008 No. 2 (Pilihan Ganda, disajikan sebagai isian)
	% Jawaban: 16
	\item (OSK 2008)
	Banyaknya faktor positif dari $5!$ adalah \dots

	% Source: OSK SMA 2007 No. 4 (Pilihan Ganda, disajikan sebagai isian)
	% Jawaban: 63
	\item (OSK 2007)
	Misalkan $H$ adalah himpunan semua faktor positif dari 2007.
	Banyaknya himpunan bagian dari $H$ yang tidak kosong adalah \dots

	% Source: OSK SMA 2012 No. 12 (Isian Singkat)
	% Jawaban: 420
	\item (OSK 2012)
	Ada berapa faktor positif dari $2^{7}3^{5}5^{3}7^{2}$
	yang merupakan kelipatan 6?

	% Source: OSK SMA 2011 No. 20 (Isian Singkat)
	% Jawaban: 378
	\item (OSK 2011)
	Ada berapa faktor positif dari $2^{7}3^{5}5^{3}7^{2}$
	yang merupakan kelipatan 10?

	% Source: OSK SMA 2010 No. 12 (Isian Singkat)
	% Jawaban: 16/81
	\item (OSK 2010)
	Misalkan $S$ menyatakan himpunan semua faktor positif dari $2010^{2}$.
	Sebuah bilangan diambil secara acak dari $S$.
	Peluang bilangan yang terambil habis dibagi 2010 adalah \dots

	% Source: OSK SMA 2017 Medan Ulang No. 5 (Isian Singkat)
	% Jawaban: 7^2016
	\item (OSK 2017)
	Suatu bilangan asli mempunyai 2017 pembagi positif yang salah satu
	di antaranya adalah 7. Bilangan yang dimaksud adalah \dots

	% Source: OSP SMA 2008 No. 1 (Isian Singkat)
	% Jawaban: 8
	\item (OSP 2008)
	Banyaknya pembagi positif dari 2008 adalah \dots

	% Source: AMC 10 2017 B (diterjemahkan, disajikan sebagai isian)
	% Jawaban: 1/19
	\item (AMC 10B 2017)
	Bilangan $21!$ memiliki lebih dari $60000$ faktor positif.
	Satu faktor dipilih secara acak.
	Peluang faktor yang terpilih ganjil adalah \dots

	% Source: AMC 10 2020 A (diterjemahkan, disajikan sebagai isian)
	% Jawaban: 23
	\item (AMC 10A 2020)
	Sebuah faktor positif dari $12!$ dipilih secara acak.
	Peluang faktor tersebut kuadrat sempurna (\textit{perfect square}) dapat ditulis $\frac{m}{n}$
	dalam bentuk paling sederhana. Nilai $m + n$ adalah \dots

	% Source: AMC 10 2019 A (diterjemahkan, disajikan sebagai isian)
	% Jawaban: 37
	\item (AMC 10A 2019)
	Ada berapa faktor positif dari $201^{9}$ yang merupakan kuadrat sempurna
	atau kubik sempurna (\textit{perfect cube}; atau keduanya)?

	% Source: AMC 12 2005 B (diterjemahkan, disajikan sebagai isian)
	% Jawaban: 2
	\item (AMC 12B 2005)
	Bilangan asli $n$ memiliki $60$ faktor positif dan $7n$ memiliki
	$80$ faktor positif. Bilangan bulat $k$ terbesar sehingga $7^{k}$
	membagi $n$ adalah \dots

	% Source: AIME 1988 No. 5 (diterjemahkan)
	% Jawaban: 634
	\item (AIME 1988)
	Peluang sebuah faktor positif dari $10^{99}$ yang dipilih secara acak
	merupakan kelipatan dari $10^{88}$ adalah $\frac{m}{n}$ dalam bentuk
	paling sederhana. Nilai $m + n$ adalah \dots

	% Source: AIME 1995 No. 6 (diterjemahkan)
	% Jawaban: 589
	\item (AIME 1995)
	Misalkan $n = 2^{31}3^{19}$. Ada berapa faktor positif dari $n^2$
	yang kurang dari $n$ tetapi tidak membagi $n$?

	% Source: OSN SMA 2004 No. 1 (Esai, disajikan sebagai isian)
	% Jawaban: 36 genap dan 12 ganjil
	\item (OSN 2004)
	Ada berapa pembagi genap dan pembagi ganjil dari $5^{6} - 1$?

	% Source: AIME 2005 I (diterjemahkan)
	% Jawaban: 25
	\item (AIME I 2005)
	Untuk bilangan asli $n$, notasi $\tau(n)$ menyatakan banyak faktor positif
	dari $n$. Definisikan $S(n) = \tau(1) + \tau(2) + \dots + \tau(n)$.
	Misalkan $a$ banyak bilangan asli $n \le 2005$ dengan $S(n)$ ganjil
	dan $b$ banyaknya dengan $S(n)$ genap. Nilai $|a - b|$ adalah \dots

	% Source: OSK SMA 2020 No. 6 (Isian Singkat)
	% Jawaban: 2
	\item (OSK 2020)
	Kuartet bilangan asli $(a, b, c, d)$ dikatakan keren jika memenuhi
	$b = a^{2} + 1$, $c = b^{2} + 1$, $d = c^{2} + 1$ dan
	$\tau(a) + \tau(b) + \tau(c) + \tau(d)$ bilangan ganjil.
	Banyaknya kuartet keren $(a, b, c, d)$ dengan $a, b, c, d < 10^{6}$
	adalah \dots\ (Catatan: $\tau(k)$ menyatakan banyaknya faktor positif dari $k$.)

	% Source: AIME 2019 I (diterjemahkan)
	% Jawaban: 540
	\item (AIME I 2019)
	Notasi $\tau(n)$ menyatakan banyak faktor positif dari $n$.
	Tentukan jumlah dari enam bilangan asli terkecil $n$ yang memenuhi
	$\tau(n) + \tau(n + 1) = 7$.

\end{enumerate}

\section{Jumlah Faktor}

\begin{enumerate}[resume]

	% Source: AMC 10 2008 A (diterjemahkan, disajikan sebagai isian)
	% Jawaban: 6
	\item (AMC 10A 2008)
	Untuk bilangan asli $n$, notasi $\langle n \rangle$ menyatakan jumlah
	semua faktor positif dari $n$ kecuali $n$ sendiri.
	Sebagai contoh, $\langle 4 \rangle = 1 + 2 = 3$ dan
	$\langle 12 \rangle = 1 + 2 + 3 + 4 + 6 = 16$.
	Nilai dari $\langle \langle \langle 6 \rangle \rangle \rangle$ adalah \dots

	% Source: AMC 10 2024 B (diterjemahkan, disajikan sebagai isian)
	% Jawaban: 6
	\item (AMC 10B 2024)
	Misalkan $N$ adalah hasil kali semua faktor positif dari $42$.
	Digit satuan dari $N$ adalah \dots

	% Source: OSK SMA 2023 No. 16 (Isian Singkat)
	% Jawaban: 27
	\item (OSK 2023)
	Misalkan $N = 2^{a} \cdot 3^{b}$ dengan $a, b$ bilangan asli.
	Jika hasil kali semua faktor positif dari $N$ adalah $24^{60}$,
	maka nilai $ab$ adalah \dots

	% Source: AMC 12 2008 B (diterjemahkan, disajikan sebagai isian)
	% Jawaban: 11
	\item (AMC 12B 2008)
	Jumlah logaritma basis $10$ dari semua faktor positif dari $10^{n}$
	adalah $792$. Nilai $n$ adalah \dots

	% Source: AIME 1986 No. 8 (diterjemahkan)
	% Jawaban: 141
	\item (AIME 1986)
	Misalkan $S$ adalah jumlah logaritma basis $10$ dari semua faktor
	positif sejati dari $1000000$ (yaitu semua faktor kecuali bilangan itu
	sendiri). Bilangan bulat terdekat ke $S$ adalah \dots

	% Source: AMC 10 2021 B (diterjemahkan, disajikan sebagai isian)
	% Jawaban: 1/14
	\item (AMC 10B 2021)
	Misalkan $N = 34 \cdot 34 \cdot 63 \cdot 270$.
	Perbandingan jumlah faktor ganjil dari $N$ terhadap jumlah faktor genap
	dari $N$ adalah \dots

	% Source: AMC 12 2021 Fall B (diterjemahkan, disajikan sebagai isian)
	% Jawaban: 1/192
	\item (AMC 12 2021 Fall B)
	Untuk bilangan asli $n$, misalkan $f(n)$ adalah hasil bagi jumlah semua
	faktor positif dari $n$ oleh $n$.
	Sebagai contoh, $f(14) = (1 + 2 + 7 + 14) \div 14 = \frac{12}{7}$.
	Nilai dari $f(768) - f(384)$ adalah \dots

	% Source: AIME 2000 I (diterjemahkan)
	% Jawaban: 248
	\item (AIME I 2000)
	Misalkan $S$ adalah jumlah semua bilangan berbentuk $\frac{a}{b}$ dengan
	$a$ dan $b$ faktor positif dari $1000$ yang relatif prima (\textit{coprime}).
	Bilangan bulat terbesar yang tidak melebihi $\frac{S}{10}$ adalah \dots

	% Source: OSK SMA 2026 No. 15 (Isian Singkat)
	% Jawaban: 2010
	\item (OSK 2026)
	Misalkan $\sigma(n)$ adalah jumlah faktor positif dari $n$.
	Jumlah semua bilangan asli $n \le 225$ sehingga $\sigma(n)$ adalah
	bilangan ganjil adalah \dots

	% Source: AIME 2021 I (diterjemahkan)
	% Jawaban: 125
	\item (AIME I 2021)
	Notasi $\sigma(a)$ menyatakan jumlah faktor positif dari $a$.
	Misalkan $n$ bilangan asli terkecil sehingga $\sigma(a^{n}) - 1$
	habis dibagi $2021$ untuk setiap bilangan asli $a$.
	Jumlah faktor prima dalam faktorisasi prima dari $n$ adalah \dots

	% Source: Modul Pra-OSP FungsiAritmetika (juga di Khurmi Bab 3 Latihan 3.7.5)
	\item (Pra-OSP Fungsi Aritmetika)
	Buktikan bahwa untuk setiap bilangan asli $n$ berlaku
	\[ \frac{\sigma(n)}{\tau(n)} \ge \sqrt{n}. \]

\end{enumerate}

\section{Fungsi Totient}

\begin{enumerate}[resume]

	% Source: AMO 2022 Grade 9 No. 20 (diterjemahkan)
	% Jawaban: 24
	\item (AMO 2022)
	Untuk bilangan asli $n$, notasi $\phi(n)$ menyatakan banyak bilangan asli
	yang kurang dari $n$ dan relatif prima dengan $n$.
	Nilai dari $\phi(35)$ adalah \dots

	% Source: OSP SMA 2008 No. 14 (Isian Singkat)
	% Jawaban: 10
	\item (OSP 2008)
	Banyak bilangan bulat positif $n$ kurang dari 2008 yang mempunyai tepat
	$\frac{n}{2}$ bilangan kurang dari $n$ dan relatif prima terhadap $n$
	adalah \dots

	% Source: Komposisi sendiri (latihan klasik)
	% Jawaban: 7, 9, 14, 18
	\item
	Tentukan semua bilangan asli $n$ sehingga $\varphi(n) = 6$.

	% Source: Modul Pra-OSP FungsiAritmetika
	\item (Pra-OSP Fungsi Aritmetika)
	Buktikan bahwa ada tak hingga bilangan asli $n$ yang memenuhi
	$\varphi(n) = \frac{n}{3}$.

\end{enumerate}

\section{Campuran}

\begin{enumerate}[resume]

	% Source: OSK SMA 2024 No. 13 (Isian Singkat)
	% Jawaban: 90
	\item (OSK 2024)
	Untuk setiap bilangan asli $n$, misalkan $f(n)$ menyatakan faktor ganjil
	terbesar dari $n$ dan $p(n) = f(n) + f(n+1) + \dots + f(2n)$.
	Jika $p(n) = 8145$, maka nilai dari $n$ adalah \dots

	% Source: AIME 1987 No. 3 (diterjemahkan)
	% Jawaban: 182
	\item (AIME 1987)
	Faktor sejati (\textit{proper divisor}) dari bilangan asli adalah faktor positif selain $1$ dan
	bilangan itu sendiri. Bilangan asli lebih dari $1$ disebut \textit{nice}
	jika sama dengan hasil kali semua faktor sejatinya yang berbeda.
	Jumlah sepuluh bilangan \textit{nice} pertama adalah \dots

	% Source: AIME 2005 I (diterjemahkan)
	% Jawaban: 109
	\item (AIME I 2005)
	Ada berapa bilangan asli yang tepat memiliki tiga faktor sejati
	(faktor positif kecuali bilangan itu sendiri),
	yang masing-masing kurang dari $50$?

	% Source: AMC 10 2016 A (diterjemahkan, disajikan sebagai isian)
	% Jawaban: 325
	\item (AMC 10A 2016)
	Untuk suatu bilangan asli $n$, bilangan $110n^{3}$ memiliki $110$ faktor
	positif. Ada berapa faktor positif dari bilangan $81n^{4}$?

	% Source: AIME 2025 II (diterjemahkan)
	% Jawaban: 237
	\item (AIME II 2025)
	Misalkan $A$ himpunan semua faktor positif dari $2025$.
	Dipilih secara acak satu himpunan bagian $B$ dari $A$.
	Peluang $B$ merupakan himpunan tak kosong yang $KPK$ semua anggotanya
	sama dengan $2025$ adalah $\frac{m}{n}$ dalam bentuk paling sederhana.
	Nilai $m + n$ adalah \dots

	% Source: OSK SMA 2024 No. 9 (Isian Singkat)
	% Jawaban: 92
	\item (OSK 2024)
	Misalkan $k$ adalah bilangan bulat positif terkecil kelipatan 2024
	yang memiliki 28 faktor positif.
	Sisa hasil bagi $k$ oleh 100 adalah \dots

	% Source: AIME 1990 No. 5 (diterjemahkan)
	% Jawaban: 432
	\item (AIME 1990)
	Misalkan $n$ bilangan asli terkecil yang merupakan kelipatan $75$
	dan memiliki tepat $75$ faktor positif.
	Nilai dari $\frac{n}{75}$ adalah \dots

	% Source: OSP SMA 2009 No. 18 (Isian Singkat)
	% Jawaban: 2
	\item (OSP 2009)
	Misalkan $n$ bilangan asli terkecil yang mempunyai tepat 2009 faktor
	dan $n$ merupakan kelipatan 2009.
	Faktor prima terkecil dari $n$ adalah \dots

	% Source: AIME 2026 I (diterjemahkan)
	% Jawaban: 244
	\item (AIME I 2026)
	Misalkan $N$ adalah banyak faktor positif dari $17017^{17}$ yang bersisa
	$5$ jika dibagi $12$. Sisa $N$ jika dibagi $1000$ adalah \dots

\end{enumerate}

\section{Tantangan}

\begin{enumerate}[resume]

	% Source: Kanada MO 1999 (via Modul Pra-OSP FungsiAritmetika)
	% Jawaban: 1 dan 9
	\item (Kanada MO 1999)
	Tentukan semua bilangan asli $n$ yang memenuhi $n = \tau(n)^2$.

	% Source: Modul Pra-OSP FungsiAritmetika
	% Jawaban: 1, 3, 8, 10, 18, 24, 30
	\item (Pra-OSP Fungsi Aritmetika)
	Carilah semua solusi persamaan $\varphi(n) = \tau(n)$.

	% Source: OSN SMA 2010 No. 6 (Esai)
	% Jawaban: 4 atau bilangan prima
	\item (OSN 2010)
	Cari semua bilangan asli $n > 1$ sedemikian sehingga
	\[ \tau(n) + \varphi(n) = n + 1. \]
	Dalam hal ini, $\tau(n)$ menyatakan banyaknya bilangan asli yang habis
	membagi $n$, dan $\varphi(n)$ menyatakan banyaknya bilangan asli yang
	kurang dari $n$ dan relatif prima terhadap $n$.

	% Source: Filipina MO 2017 (via Modul Pra-OSP FungsiAritmetika)
	% Jawaban: kesamaan tercapai jika dan hanya jika n = 1
	\item (Filipina MO 2017)
	Buktikan bahwa
	\[ \frac{1}{\sigma(n)} + \frac{1}{\varphi(n)} \ge \frac{2}{n} \]
	dan tentukan kapan kesamaan tercapai.

	% Source: Turki MO 2022 (via Modul Pra-OSP FungsiAritmetika)
	\item (Turki MO 2022)
	Buktikan bahwa untuk setiap bilangan asli $k, n$ yang memenuhi
	$k \ge n!$ berlaku $\varphi(k) \ge (n - 1)!$.

	% Source: Irlandia MO 1998 (via Khurmi Bab 3 Latihan 3.7.10)
	% Jawaban: 1998 dan 3834
	\item (Irlandia MO 1998)
	Tentukan semua bilangan asli $d$ yang memiliki tepat $16$ faktor positif
	$d_1 < d_2 < \dots < d_{16}$ sehingga $d_6 = 18$ dan $d_9 - d_8 = 17$.

\end{enumerate}

\end{document}
